\chapter{System Implementation and Methodology}
\label{chap:implementation}

\section{Development Methodology}
\label{sec:methodology}

Development moved in vertical slices: pose extraction on a single still image first, then on a video, then shot detection, then the metrics layer, then the API on top of that, then the mobile screen against the API. Frontend and backend work ran in parallel once the API contract was locked---\lstinline|src/types/contracts.ts| on the TypeScript side mirrored \lstinline|backend/contracts/*.py| on the Python side, so both tracks could move independently without drifting. Contract mismatches surfaced as TypeScript compilation errors rather than silent runtime failures.

After the core feature set stabilised, a structured production-readiness audit was conducted across the full codebase. The audit flagged 55 issues across seven severity categories. Fifty were resolved before submission; the remaining five are architectural decisions---primarily around which routes require authentication---documented with explicit remediation plans.

\section{Frontend Implementation}
\label{sec:frontend_impl}

\subsection{Application Entry Point}

\lstinline|App.tsx| wraps the application in three context providers (\lstinline|AuthProvider|, \lstinline|HistoryProvider|, \lstinline|ProfileProvider|) and delegates rendering to \lstinline|AppNavigator|, which handles the authentication gate: unauthenticated users see \lstinline|LoginScreen|; authenticated users see the main bottom-tab navigator.

\subsection{MVPAnalysisScreen}

\lstinline|MVPAnalysisScreen| is the central user-facing feature. It manages five distinct UI states: initial (file selection), uploading (progress indicator), analysing (animated pulse with rotating status labels), result (score ring and metrics), and error (descriptive message). The screen implements the client-side polling loop shown in Listing~\ref{lst:polling}.

\noindent\begin{minipage}{\linewidth}
\begin{lstlisting}[language=TypeScript,
                   caption={Client-side polling loop (simplified from \texttt{MVPAnalysisScreen.tsx}).},
                   label={lst:polling}]
const pollJobResult = async (jobId: string) => {
  while (true) {
    const result = await apiService.getAnalysisResult(jobId);
    if (result.status === 'completed') {
      setAnalysisResult(result);
      await persistAnalysis(result); // POST /api/analysis/complete
      break;
    }
    if (result.status === 'failed') {
      setError(result.error ?? 'Analysis failed');
      break;
    }
    await sleep(2000); // poll every 2 seconds
  }
};
\end{lstlisting}
\end{minipage}

The persist step applies retry logic---if the initial call fails due to a network flap, it retries up to three times before surfacing an error to the user.

\subsection{AuthContext: PKCE OAuth Flow}

\lstinline|AuthContext| implements the full PKCE OAuth flow for Google sign-in. The PKCE code verifier is generated locally using \lstinline|expo-crypto|, and the authorisation code is extracted from the deep link redirect URI (\lstinline|shootrz://...?code=...|) using regex parsing as a fallback when the standard URL API fails on certain Android device configurations.

\subsection{Design System}

A token-based design system (\lstinline|src/theme/tokens.ts|) defines the full colour palette, type scale, spacing units, and motion curves as named constants. Score tier colours---\textit{elite} (gold), \textit{great} (green), \textit{good} (blue), \textit{fair} (orange), \textit{poor} (red)---are defined as named tokens rather than inline hex values, ensuring consistent application across all components and preventing undetected token name typos via a dedicated contract test suite.

\section{Backend Implementation}
\label{sec:backend_impl}

\subsection{App Factory}

The FastAPI application is created inside an \lstinline|@asynccontextmanager| lifespan function that enforces \lstinline|multiprocessing.set_start_method("spawn")| on non-Windows platforms. This is required because MediaPipe initialises process-level state that is not fork-safe; without this setting, pipeline processes spawned on Linux will silently hang or produce corrupt pose data. JSON logging is initialised with \lstinline|logging.config.dictConfig| before any routers are imported to guarantee that all startup messages are captured in structured format.

\subsection{MVPPipeline}

The pipeline (\lstinline|backend/mvp/core/pipeline.py|) is implemented as an orchestrator class that calls six stage functions in sequence, passing a shared \lstinline|PipelineState| dictionary through each stage. Each stage writes its outputs (DataFrames, JSON dictionaries, file paths) into the state, and the next stage reads from it. This design allows any individual stage to be independently tested with synthetic state inputs. The key state fields are shown in Listing~\ref{lst:pipeline_state}.

\noindent\begin{minipage}{\linewidth}
\begin{lstlisting}[language=Python,
                   caption={Schematic \texttt{PipelineState} dictionary passed between pipeline stages.},
                   label={lst:pipeline_state}]
PipelineState = {
    "frames":         List[np.ndarray],      # Stage 1 output
    "frame_mapping":  pd.DataFrame,          # Stage 1 output
    "metadata":       VideoMetadata,         # Stage 1 output
    "pose_keypoints": pd.DataFrame,          # Stage 2 output
    "smoothed_pose":  pd.DataFrame,          # Stage 3 output
    "shot_window":    ShotWindow,            # Stage 4 output
    "angles":         pd.DataFrame,          # Stage 5 output
    "metrics":        List[MVPMetric],       # Stage 6 output
    "overall_score":  float,                 # Stage 6 output
}
\end{lstlisting}
\end{minipage}

\subsection{Signal Smoothing}

Savitzky--Golay filtering (window length 5, polynomial order 2) is applied independently to each joint's $x$, $y$, and $z$ coordinate series. This filter was chosen because it preserves the shape of sharp features---the wrist-height peak at release, the knee-angle minimum at crouch---better than a simple moving average, which would underestimate the true extrema. Gaps of up to 10 consecutive missing frames are interpolated using linear interpolation before smoothing is applied.

\subsection{LLM Job Service}

The enrichment step in \lstinline|mvp_job_service.py| wraps the Gemini API call in \lstinline|asyncio.wait_for| with a 10-second timeout. If the call completes successfully, the AI-scored overall value and feedback are merged into the pipeline result payload. If it times out or raises an exception, the payload is returned with \lstinline|ai_scored=false| and the rule-based score is used. This design ensures that LLM latency spikes never delay analysis results beyond the pipeline's own runtime.

\subsection{Supabase DB Layer}

All database interactions are centralised in a \lstinline|SupabaseDB| singleton (\lstinline|backend/storage/db.py|) that uses the service-role key, bypassing RLS for backend writes (appropriate for a trusted server process), while ensuring that the anon key is used only for client-facing JWT verification. Methods follow an upsert-or-insert pattern to handle duplicate calls arising from the client's retry logic.

\subsection{Rule-Based Feedback Engine}

A rule-based fallback engine (\lstinline|backend/feedback/|) generates coaching cues without LLM involvement. Rules are defined as threshold comparisons against the computed metric values, producing templated strings. For example, when \lstinline|elbow_angle < good_range_lower|, the engine produces: ``Your elbow is not fully extending at release---try finishing with your arm straight toward the basket.'' These rules serve as the fallback when Gemini is unavailable and as a validation baseline for LLM output quality.

\subsection{Chat Context Builder}

Context injection for the chat endpoint calls \lstinline|get_coach_context|, a Supabase RPC defined in \lstinline|backend/chat/context_builder.py|, which returns the user's recent analyses and metrics in a single database round trip. Previously, this required four separate queries (analyses, metrics, profile, streak). The context is formatted as a structured prompt prefix that Gemini receives before the user's message, giving the LLM specific awareness of the user's history, weak areas, and skill level.

\section{AI Pipeline Implementation}
\label{sec:ai_pipeline}

\subsection{Preprocessing}

Raw video frames are decoded by OpenCV and converted from BGR to RGB colour order. A stride-based frame sampler limits processing to at most 300 frames per video:

\begin{equation}
  \text{stride} = \left\lceil \frac{\text{total\_frames}}{300} \right\rceil
\end{equation}

This caps memory usage at a predictable level while preserving the frame-to-timestamp mapping for accurate metric plotting.

\subsection{Pose Estimation}

MediaPipe Pose is initialised once per process (not per video) to amortise the model loading cost. Each frame is passed to \lstinline|mediapipe.solutions.pose.Pose.process()|, which returns 33 normalised landmarks. Results are written to a DataFrame with columns \lstinline|[frame_id, joint, x, y, z, confidence]| where joint names follow the SHOOTRZ basketball-specific naming convention (e.g., \lstinline|right_elbow|, \lstinline|left_knee|).

\subsection{Feature Extraction}

Angle features are computed by the \lstinline|AngleComputer| class in Stage~5. For a triplet of joints A, B, C (angle measured at B):

\begin{equation}
  \theta = \arccos\!\left(\operatorname{clip}\!\left(\frac{(\mathbf{A}-\mathbf{B})\cdot(\mathbf{C}-\mathbf{B})}{\|\mathbf{A}-\mathbf{B}\|\;\|\mathbf{C}-\mathbf{B}\|},\; -1,\; 1\right)\right)
\end{equation}

For elbow extension, $\theta = 180^{\circ}$ indicates a fully extended arm.

\subsection{Output Artefacts}

The pipeline writes six artefact files to a per-run output directory (\lstinline|backend/outputs/\{run\_id\}/|): \lstinline|config_used.yaml|, \lstinline|video_metadata.json|, \lstinline|frame_mapping.csv|, \lstinline|pose_keypoints.csv|, \lstinline|angles.csv|, and \lstinline|run_metadata.json|. The overlay video (\lstinline|overlay_video.mp4|) is generated by \lstinline|video_annotator.py| using OpenCV to draw the MediaPipe skeleton, metric verdicts, and score overlay onto each original frame.

\subsection{Score Computation: Declared but Unused Dependencies}

Scores are derived deterministically from computed angle values using the normative range functions; they are not predicted by a trained classifier. The LightGBM~\cite{Ke2017} library is declared in \lstinline|requirements.txt| for a planned shot success prediction model, and PyTorch is declared for a Phase~3 ML refinement pipeline. Neither is currently instantiated in any active inference path; they represent design intent for future capability rather than implemented features.

\section{Technologies Used}
\label{sec:technologies}

Table~\ref{tab:tech_stack} lists the core technologies used in SHOOTRZ with the version deployed and the engineering rationale for each choice.

\begin{longtable}{>{\raggedright\arraybackslash}p{0.19\linewidth}>{\raggedright\arraybackslash}p{0.11\linewidth}>{\raggedright\arraybackslash}p{0.58\linewidth}}
  \caption{Technology stack with engineering rationale.}
  \label{tab:tech_stack} \\
  \toprule
  \textbf{Technology} & \textbf{Version} & \textbf{Rationale} \\
  \midrule
  \endfirsthead
  \toprule
  \textbf{Technology} & \textbf{Version} & \textbf{Rationale} \\
  \midrule
  \endhead
  \bottomrule
  \endfoot
  MediaPipe Pose & 0.10.14 & State-of-the-art CPU-viable monocular pose estimation; no GPU required \\
  FastAPI & 0.110+ & Async-first Python HTTP framework with native Pydantic v2 integration \\
  Pydantic v2 & 2.6+ & Strict data validation at API boundaries; TypeScript-compatible schema export \\
  Expo / React Native & 54 / 0.81.5 & Cross-platform mobile with managed EAS build pipeline \\
  Supabase & --- & Managed PostgreSQL with auth, storage, and RLS in a single platform \\
  Gemini 2.5 Flash & --- & Best cost-to-capability ratio for structured JSON output generation \\
  FAISS & 1.7.4 & Sub-millisecond similarity search for small-to-medium drill embeddings \\
  MABWiser (LinUCB) & 0.4+ & Contextual bandit implementation with a clean numpy backend \\
  SciPy & 1.11+ & \lstinline|find_peaks| and \lstinline|savgol_filter| for signal processing \\
  YOLOv8n & 8.0+ & Lightweight object detector suitable for CPU ball tracking \\
  slowapi & 0.1.9 & FastAPI-native rate limiting via decorator API \\
\end{longtable}

\section{Implementation Challenges}
\label{sec:challenges}

\subsection{Shooting Side Detection}

The initial implementation hardcoded right-side landmark indices in several inference modules. Correcting this required threading a \lstinline|shooting_side| parameter through the entire pipeline---VideoLoader $\rightarrow$ PoseEstimation $\rightarrow$ ShotDetection $\rightarrow$ AngleComputation---and updating all joint name lookups to use the \lstinline|{shooting_side}_{joint}| convention. Left-handed players had been receiving systematically wrong elbow and wrist angle values before this fix.

\subsection{Geometric Mean Collapse}

Early testing revealed that the geometric mean produced near-zero overall scores when any one metric was missing or had very low confidence, even if the rest of the shot was excellent. The fix involved separating the confidence-gating logic (which discards unreliable frames from metric selection) from the scoring logic (which only uses metrics with valid selected frames), and implementing the Gemini AI score override as a practical correction for edge cases where a single low-confidence axis collapses an otherwise valid score.

\subsection{MediaPipe Fork Safety}

FastAPI's default worker model on Linux uses \lstinline|fork()|, which is incompatible with MediaPipe's process-level C++ runtime initialisation. Diagnosing this required understanding the interaction between \lstinline|os.fork()|, MediaPipe's shared native state, and Python's \lstinline|multiprocessing| module. \begin{sloppypar}
Enforcing \lstinline|set_start_method("spawn")| in the lifespan context manager fixed this: a one-line change that was non-obvious to diagnose but critical to production stability.
\end{sloppypar}

\subsection{SSE Multi-Line Parsing}

The frontend chat service initially parsed server-sent event data assuming each \lstinline|data:| prefix would contain a complete JSON object. Gemini's streaming output sometimes spans multiple lines within a single event, causing JSON parse errors at chunk boundaries. The fix concatenates successive \lstinline|data:| lines before attempting JSON parsing.

\subsection{PKCE Deep Link Parsing}

On certain Android devices, the OAuth redirect URI is delivered to the application before the JavaScript runtime has fully initialised. This produced intermittent failures where the PKCE authorisation code was silently lost. The fix added a regex-based fallback parser alongside the standard URL API, ensuring the code is captured regardless of initialisation timing.
